
\noindent A \textit{Pisot} (or Pisot--Vijayaraghavan) number is a real algebraic integer \(\nu > 1\) whose conjugates all lie within the unit circle in the complex plane. For example, the golden mean \(\varphi = (1+\sqrt{5})/2\) is a Pisot number because its conjugate \(\varphi' = (1-\sqrt{5})/2 \approx -0.618\) has modulus less than \(1\). Pisot numbers play a significant role in the study of substitution tilings. Specifically, if the substitution factor of a tiling (from a fairly general class of substitution tilings with fractagonal tiles) is a Pisot number, then the tiling exhibits \textit{finite local complexity} (FLC). This result was proven by Frank and Robinson~\cite{FrankRob05}, with a special case given by Frettl\"oh~\cite{Fre02}. A tiling is said to have FLC if it contains only finitely many distinct two-tile patches, up to rigid motions. (Note that some authors define FLC considering only translations.) A tiling that does not exhibit FLC, meaning it has infinitely many two-tile patches, is said to have \textit{infinite local complexity} (ILC).

The converse of the statement “Pisot substitution factor implies FLC” is not generally true. There exist substitution tilings with non-Pisot substitution factors, yet still exhibit FLC. Examples include the chiral Lan\c{c}on--Billard tiling~\cite{LanBill88, BaaGri13} and certain octagonal tilings recently found by Say-awen and Coates~\cite{SayCo}. However, when the additional condition is imposed that some adjacent tiles do not meet full-edge to full-edge, the tiling can be shown to have ILC~\cite{Danz02, FrankRob05}.

In this work, we contribute to the study of substitution tilings by presenting three families of substitution tilings with ILC. We utilize Danzer’s algorithm~\cite{Danz02}, which applies when the substitution factor is non-Pisot and the tiling contains a vertex in a \textit{misfit situation}. A vertex \(\mathcal{V}\) of a tiling \(\mathcal{T}\) is in a misfit situation if there exists a tile \(T \in \mathcal{T}\) such that \(\mathcal{V}\) lies in \(T\) but is not one of its vertices. Equivalently, \(\mathcal{T}\) contains two adjacent tiles that do not meet full-edge to full-edge.

The tilings in this work are constructed using substitution rules whose prototiles are isosceles triangles with interior angles that are multiples of \(\pi/n\), and whose substitution factor is the length of the longest diagonal of a regular \(n\)-gon with unit side length. A more general construction of tilings with this substitution factor and a broader class of prototiles was given by Nischke and Danzer~\cite{NisDanz96}. Similar constructions were presented by Pautze~\cite{Pau17} for \(n = 7, 11\) and by Hibma~\cite{Hib15} for odd \(n\) with \(5 \leq n \leq 15\).

In addition to exhibiting ILC, a distinctive feature of each substitution described here is that its tiling space contains a tiling with \(n\)-fold rotational symmetry. This symmetry distinguishes my examples from most of those in~\cite{NisDanz96} and~\cite{Hib15}. To the best of my knowledge, these are the first examples of substitution tilings with \(n\)-fold rotational symmetry for these specific values of \(n\).
	
